%%%%%%%%%%%%%%%%%%%%%%%%%%%%%%%%%%%%%%%%%%%%%%%%%%%%%%%%%%%%%%%%%%%%%%%%%%%
\section{Computational Specification of Orthogonal Projection Estimation}
\label{app:computational_specification}
%%%%%%%%%%%%%%%%%%%%%%%%%%%%%%%%%%%%%%%%%%%%%%%%%%%%%%%%%%%%%%%%%%%%%%%%%%%

This appendix provides the formal computational specification underlying the
orthogonal projection estimators developed in Sections~\ref{sec:projection}--
\ref{sec:algorithm}. Its purpose is to translate the statistical objects of the
main text into a model-independent collection of computational inputs,
operations, outputs, and diagnostics. The specification is not tied to a
particular programming language. It instead defines the mathematical contract
that an implementation must satisfy. The \texttt{OrthoGMM} package is one
software realization of this contract.

The appendix has six objectives. First, it separates the statistical target
from the numerical procedure used to compute it. Second, it identifies the
unit-level moment contributions required for projection and covariance
estimation. Third, it formalizes the tractable, demanding, derivative, and
reconstruction operators supplied by a model. Fourth, it specifies the
localization, projection, correction, and inference workflow. Fifth, it defines
numerical safeguards and computational diagnostics. Finally, it shows how the
Berry--Levinsohn--Pakes application instantiates the general specification.

Throughout, the statistical definitions and efficiency results are those of
Sections~\ref{sec:projection} and~\ref{sec:gmm}. The present appendix does not
introduce an alternative estimator. It fixes a single feasible implementation
of the projected efficient estimator and records the information required to
reproduce it.

%%%%%%%%%%%%%%%%%%%%%%%%%%%%%%%%%%%%%%%%%%%%%%%%%%%%%%%%%%%%%%%%%%%%%%%%%%%
\subsection{Computational Design Principles}
\label{app:design_principles}
%%%%%%%%%%%%%%%%%%%%%%%%%%%%%%%%%%%%%%%%%%%%%%%%%%%%%%%%%%%%%%%%%%%%%%%%%%%

\paragraph{Separation of statistical and computational structure.}
Let $\theta\in\Theta\subset\mathbb R^p$ denote the parameter of interest. The
statistical model determines the valid estimating equations, the efficient
influence function, and the asymptotic variance bound. The computational
structure determines how these objects are evaluated. The specification keeps
these layers distinct:
\[
\text{statistical model}
\quad\longrightarrow\quad
\text{estimating equations and efficiency target},
\]
whereas
\[
\text{computational implementation}
\quad\longrightarrow\quad
\text{operators, numerical tolerances, and evaluation costs}.
\]
Two implementations are therefore statistically equivalent when they evaluate
the same estimating equations to the accuracy required by the theory, even if
they use different numerical methods, programming languages, or hardware.

\paragraph{Computational partition.}
For each elementary statistical unit $i$, let
\begin{equation}
 m_i(\theta)
 =
 \begin{pmatrix}
 g_i(\theta)\\
 h_i(\theta)
 \end{pmatrix},
 \qquad
 g_i(\theta)\in\mathbb R^{q_g},
 \quad
 h_i(\theta)\in\mathbb R^{q_h},
 \label{eq:app_moment_partition}
\end{equation}
where $g_i$ is computationally tractable and $h_i$ is computationally
demanding. The partition is defined by evaluation cost rather than statistical
importance. Both components are valid estimating equations at the true
parameter:
\begin{equation}
 E[g_i(\theta_0)]=0,
 \qquad
 E[h_i(\theta_0)]=0.
 \label{eq:app_valid_moments}
\end{equation}
The tractable subsystem must be sufficiently informative to deliver a
root-$n$ consistent localizer. The demanding subsystem contributes the
incremental information required for full efficiency.

\paragraph{Unit-level contributions.}
Projection and covariance estimation require the collection
$\{g_i(\theta),h_i(\theta)\}_{i=1}^n$, not only their sample averages. An
implementation must therefore expose moment contributions at the elementary
statistical unit. Depending on the application, this unit may be an
observation, cluster, market, household, firm, panel unit, or another sampling
block. The definition of the unit determines the covariance convention and
must be fixed before estimation.

\paragraph{Local use of demanding computation.}
The demanding operator is not placed inside a global optimizer. It is evaluated
only after the tractable subsystem has localized the parameter. The resulting
workflow is
\[
\text{global tractable localization}
\;\longrightarrow\;
\text{local demanding evaluation}
\;\longrightarrow\;
\text{projected efficient correction}.
\]
This is the defining computational restriction of the implementation. A model
may use several demanding evaluations to obtain derivatives or sensitivity
checks, but those evaluations must remain local to the preliminary estimator.

\paragraph{Reproducible numerical state.}
Simulation-based and iterative models must make their numerical state explicit.
This includes random seeds, simulation draws, quadrature nodes, contraction
tolerances, solver tolerances, warm starts, and regularization parameters.
Common random numbers should be used across nearby parameter evaluations unless
a different design is explicitly justified. The numerical state forms part of
the estimator specification because changing it can change finite-sample moment
values and derivatives.

\paragraph{Transparent computational accounting.}
Wall-clock time alone is not an adequate measure of computational performance.
An implementation must separately count tractable evaluations, demanding
evaluations, derivative evaluations, inner-solver iterations, and final
reconstruction calls. These operation-specific counts permit comparisons that
remain meaningful across hardware and software environments.

%%%%%%%%%%%%%%%%%%%%%%%%%%%%%%%%%%%%%%%%%%%%%%%%%%%%%%%%%%%%%%%%%%%%%%%%%%%
\subsection{Statistical Objects and Computational Operators}
\label{app:statistical_objects_operators}
%%%%%%%%%%%%%%%%%%%%%%%%%%%%%%%%%%%%%%%%%%%%%%%%%%%%%%%%%%%%%%%%%%%%%%%%%%%

\subsubsection{Population projection objects}

Suppressing evaluation at $\theta_0$ where no ambiguity arises, define
\begin{align}
 \Omega_{gg} &= E[g_i g_i'],
 &
 \Omega_{gh} &= E[g_i h_i'],
 \nonumber\\
 \Omega_{hg} &= E[h_i g_i'],
 &
 \Omega_{hh} &= E[h_i h_i'].
 \label{eq:app_covariance_blocks}
\end{align}
The projection coefficient of the demanding moments on the tractable moments is
\begin{equation}
 B_0
 =
 \Omega_{hg}\Omega_{gg}^{-1}.
 \label{eq:app_projection_coefficient}
\end{equation}
The orthogonal demanding residual is
\begin{equation}
 \nu_i
 =
 h_i-B_0g_i,
 \qquad
 E[\nu_i g_i']=0,
 \label{eq:app_orthogonal_residual}
\end{equation}
and its covariance matrix is
\begin{equation}
 S_0
 =
 E[\nu_i\nu_i']
 =
 \Omega_{hh}
 -
 \Omega_{hg}\Omega_{gg}^{-1}\Omega_{gh}.
 \label{eq:app_schur_complement}
\end{equation}

Let
\begin{equation}
 G_0
 =
 E\!\left[
 \frac{\partial g_i(\theta_0)}{\partial\theta'}
 \right],
 \qquad
 H_0
 =
 E\!\left[
 \frac{\partial h_i(\theta_0)}{\partial\theta'}
 \right].
 \label{eq:app_jacobians}
\end{equation}
The projected demanding Jacobian is
\begin{equation}
 R_0
 =
 H_0-B_0G_0.
 \label{eq:app_projected_jacobian}
\end{equation}
The projected information matrix is
\begin{equation}
 J_0
 =
 G_0'\Omega_{gg}^{-1}G_0
 +
 R_0'S_0^{-1}R_0.
 \label{eq:app_information_matrix}
\end{equation}
Under the conditions of Section~\ref{sec:gmm}, this matrix is algebraically
equivalent to the efficient GMM information matrix based on the full moment
vector.

\subsubsection{Feasible sample objects}

Let $\widetilde\theta$ denote the tractable preliminary estimator. Define the
unit-level evaluations
\[
 \widetilde g_i=g_i(\widetilde\theta),
 \qquad
 \widetilde h_i=h_i(\widetilde\theta),
\]
and the corresponding sample averages
\[
 \overline g
 =
 \frac{1}{n}\sum_{i=1}^n\widetilde g_i,
 \qquad
 \overline h
 =
 \frac{1}{n}\sum_{i=1}^n\widetilde h_i.
\]
For independent elementary units, the baseline covariance estimators are
\begin{align}
 \widehat\Omega_{gg}
 &=
 \frac{1}{n}\sum_{i=1}^n
 (\widetilde g_i-\overline g)
 (\widetilde g_i-\overline g)',
 \label{eq:app_sample_omega_gg}\\
 \widehat\Omega_{hg}
 &=
 \frac{1}{n}\sum_{i=1}^n
 (\widetilde h_i-\overline h)
 (\widetilde g_i-\overline g)',
 \label{eq:app_sample_omega_hg}\\
 \widehat\Omega_{hh}
 &=
 \frac{1}{n}\sum_{i=1}^n
 (\widetilde h_i-\overline h)
 (\widetilde h_i-\overline h)'.
 \label{eq:app_sample_omega_hh}
\end{align}
When observations are dependent within clusters, these expressions are replaced
by covariance estimators based on cluster-level sums. HAC or spatial covariance
estimators may be used when required by the sampling design. The covariance
convention must be applied consistently to all blocks.

The feasible projection coefficient and residuals are
\begin{equation}
 \widehat B
 =
 \widehat\Omega_{hg}
 \widehat\Omega_{gg}^{-1},
 \qquad
 \widehat\nu_i
 =
 \widetilde h_i-
 \widehat B\widetilde g_i.
 \label{eq:app_feasible_projection}
\end{equation}
Their covariance is estimated by
\begin{equation}
 \widehat S
 =
 \frac{1}{n}\sum_{i=1}^n
 (\widehat\nu_i-\overline\nu)
 (\widehat\nu_i-\overline\nu)',
 \qquad
 \overline\nu
 =
 \frac{1}{n}\sum_{i=1}^n\widehat\nu_i.
 \label{eq:app_feasible_S}
\end{equation}
Let $\widehat G$ and $\widehat H$ denote feasible estimates of the two
Jacobians at $\widetilde\theta$. Then
\begin{equation}
 \widehat R
 =
 \widehat H-
 \widehat B\widehat G,
 \qquad
 \widehat J
 =
 \widehat G'\widehat\Omega_{gg}^{-1}\widehat G
 +
 \widehat R'\widehat S^{-1}\widehat R.
 \label{eq:app_feasible_R_J}
\end{equation}
The feasible projected score is
\begin{equation}
 \widehat\psi_n
 =
 \widehat G'\widehat\Omega_{gg}^{-1}\overline g
 +
 \widehat R'\widehat S^{-1}\overline\nu,
 \label{eq:app_projected_score}
\end{equation}
where $\overline\nu=n^{-1}\sum_i\widehat\nu_i$. The canonical one-step
estimator implemented by the package is
\begin{equation}
 \widehat\theta_{\mathrm{SEIP}}
 =
 \widetilde\theta
 -
 \widehat J^{-1}\widehat\psi_n.
 \label{eq:app_seip_estimator}
\end{equation}
A damped version may be used for finite-sample numerical stability:
\begin{equation}
 \widehat\theta_{\mathrm{SEIP}}(\alpha)
 =
 \widetilde\theta
 -
 \alpha\widehat J^{-1}\widehat\psi_n,
 \qquad
 \alpha\in(0,1].
 \label{eq:app_damped_update}
\end{equation}
The undamped update $\alpha=1$ is the default. Any departure from it must be
reported.

\subsubsection{Computational operators}

A model implementation supplies the following four operators.

\paragraph{Tractable moment operator.}
The operator
\begin{equation}
 \mathcal G:
 \Theta\times\mathcal Z
 \longrightarrow
 \mathbb R^{n\times q_g}
 \label{eq:app_G_operator}
\end{equation}
maps a parameter value and fixed numerical state $\mathcal Z$ into the matrix
whose $i$th row is $g_i(\theta)'$. It must support repeated global evaluation
and return unit-level contributions in a stable order.

\paragraph{Demanding moment operator.}
The operator
\begin{equation}
 \mathcal H:
 \Theta\times\mathcal Z
 \longrightarrow
 \mathbb R^{n\times q_h}
 \label{eq:app_H_operator}
\end{equation}
returns the matrix whose $i$th row is $h_i(\theta)'$. Evaluation may invoke
simulation, integration, fixed-point iteration, equilibrium solution, or nested
optimization. It must also return sufficient metadata to account for the
numerical work performed.

\paragraph{Derivative operator.}
The operator
\begin{equation}
 \mathcal D:
 \Theta\times\mathcal Z
 \longrightarrow
 \mathbb R^{q_g\times p}
 \times
 \mathbb R^{q_h\times p}
 \label{eq:app_D_operator}
\end{equation}
returns estimates of $G(\theta)$ and $H(\theta)$. The implementation must
record whether derivatives are analytical, obtained by automatic
differentiation, or approximated numerically. If numerical derivatives are
used, the step rule and any parameter scaling must be stored with the results.

\paragraph{Reconstruction operator.}
The operator
\begin{equation}
 \mathcal R:
 \Theta\times\mathcal Z
 \longrightarrow
 \mathcal Q
 \label{eq:app_R_operator}
\end{equation}
constructs model-specific quantities $\mathcal Q$ at the final estimator.
Examples include latent variables, elasticities, marginal effects,
counterfactual equilibria, welfare measures, or efficiency scores. The
reconstruction operator is not part of the estimating correction unless the
requested quantity is itself required to evaluate the moments. Its calls are
therefore counted separately.

\paragraph{Operator output contract.}
Each moment operator must return at least
\begin{enumerate}
 \item the unit-level moment matrix;
 \item the aggregate sample moment;
 \item dimensions and unit identifiers;
 \item success or failure status;
 \item numerical tolerances actually achieved;
 \item operation counts and elapsed time; and
 \item the numerical state required to reproduce the evaluation.
\end{enumerate}
A failed inner solution must never be silently treated as a valid moment
evaluation.

%%%%%%%%%%%%%%%%%%%%%%%%%%%%%%%%%%%%%%%%%%%%%%%%%%%%%%%%%%%%%%%%%%%%%%%%%%%
\subsection{Computational Workflow}
\label{app:computational_workflow}
%%%%%%%%%%%%%%%%%%%%%%%%%%%%%%%%%%%%%%%%%%%%%%%%%%%%%%%%%%%%%%%%%%%%%%%%%%%

The canonical implementation contains six stages. The stages are fixed even
when model-specific numerical methods differ.

\subsubsection{Stage 1: tractable localization}

Construct $\widetilde\theta$ using only the tractable subsystem. For example,
with a positive definite tractable weighting matrix $W_g$,
\begin{equation}
 \widetilde\theta
 =
 \arg\min_{\theta\in\Theta}
 \overline g_n(\theta)'W_g\overline g_n(\theta),
 \qquad
 \overline g_n(\theta)
 =
 \frac{1}{n}\sum_{i=1}^n g_i(\theta).
 \label{eq:app_localizer}
\end{equation}
The optimization method, parameter bounds, initial values, convergence
criterion, and attained criterion value must be retained. The package should
not proceed automatically when localization fails or returns a parameter on an
unintended boundary.

\subsubsection{Stage 2: local demanding evaluation}

Evaluate $\mathcal H$ at $\widetilde\theta$. When $\widehat H$ is obtained by
finite differences, additional evaluations are made at local perturbations
\[
 \widetilde\theta\pm\delta_j e_j,
 \qquad j=1,\ldots,p.
\]
All such perturbations must remain inside the admissible parameter space. The
same simulation draws or quadrature nodes should be used across perturbations
unless the user explicitly selects independent numerical states.

\subsubsection{Stage 3: projection construction}

Using unit-level outputs from $\mathcal G$ and $\mathcal H$, estimate
$\widehat\Omega_{gg}$, $\widehat\Omega_{hg}$, $\widehat B$, and
$\widehat S$. Compute the orthogonality residual
\begin{equation}
 \widehat\Delta_{g\nu}
 =
 \frac{1}{n}\sum_{i=1}^n
 (\widetilde g_i-\overline g)
 (\widehat\nu_i-\overline\nu)'.
 \label{eq:app_orthogonality_residual}
\end{equation}
The norm of this matrix is a required diagnostic. Large values indicate an
inconsistent unit definition, covariance convention, numerical failure, or
insufficient regularization accuracy.

\subsubsection{Stage 4: projected efficient correction}

Estimate $\widehat G$, $\widehat H$, $\widehat R$, and $\widehat J$, then
compute~\eqref{eq:app_projected_score} and~\eqref{eq:app_seip_estimator}. Linear
systems should be solved directly rather than by explicitly forming matrix
inverses whenever practical. The implementation must retain the update vector
\begin{equation}
 \widehat d
 =
 -\widehat J^{-1}\widehat\psi_n
 \label{eq:app_update_vector}
\end{equation}
and its scaled norm. A large update relative to the scale of
$\widetilde\theta$ is evidence that the tractable localizer or the selected
computational partition is inadequate.

\subsubsection{Stage 5: inference}

Under efficient weighting and the conditions of Section~\ref{sec:gmm}, the
baseline covariance estimator is
\begin{equation}
 \widehat{\operatorname{Var}}
 (\widehat\theta_{\mathrm{SEIP}})
 =
 \frac{1}{n}\widehat J^{-1}.
 \label{eq:app_variance_estimator}
\end{equation}
When non-efficient weights, clustered covariance estimators, or estimated
numerical approximations require a sandwich form, the implementation must use
the corresponding influence-function covariance rather than
mechanically applying~\eqref{eq:app_variance_estimator}. Standard errors,
confidence intervals, and covariance matrices must all identify the covariance
convention used.

\subsubsection{Stage 6: final reconstruction}

Evaluate $\mathcal R$ at $\widehat\theta_{\mathrm{SEIP}}$. When a high-fidelity
reconstruction is requested, it is counted separately from the demanding
evaluations used in projection and differentiation. This distinction prevents
the cost of reporting economic quantities from being conflated with the cost of
estimation.

\subsubsection{Canonical pseudocode}

The complete workflow is summarized by the following implementation-independent
pseudocode.

\begin{enumerate}
 \item Fix the data, elementary statistical unit, covariance convention, and
 numerical state $\mathcal Z$.
 \item Obtain $\widetilde\theta$ by optimizing the tractable criterion.
 \item Evaluate unit-level tractable and demanding moments at
 $\widetilde\theta$.
 \item Evaluate or approximate the two Jacobians locally.
 \item Estimate covariance blocks, the projection coefficient, orthogonal
 residuals, the residual covariance, and the projected Jacobian.
 \item Form the projected score and solve for the one-step correction.
 \item Compute inference using the selected covariance convention.
 \item Reconstruct requested structural quantities at the final estimator.
 \item Return estimates, inference, diagnostics, numerical state, and complete
 evaluation counts.
\end{enumerate}

%%%%%%%%%%%%%%%%%%%%%%%%%%%%%%%%%%%%%%%%%%%%%%%%%%%%%%%%%%%%%%%%%%%%%%%%%%%
\subsection{Numerical Specification and Diagnostics}
\label{app:numerical_specification}
%%%%%%%%%%%%%%%%%%%%%%%%%%%%%%%%%%%%%%%%%%%%%%%%%%%%%%%%%%%%%%%%%%%%%%%%%%%

\subsubsection{Derivative evaluation}

Analytical derivatives are preferred when available. Automatic differentiation
is appropriate when the underlying numerical routine is differentiable and the
resulting derivatives have been validated. Otherwise, central differences are
the default:
\begin{equation}
 \frac{\partial \overline h_n(\theta)}{\partial\theta_j}
 \approx
 \frac{
 \overline h_n(\theta+\delta_j e_j)
 -
 \overline h_n(\theta-\delta_j e_j)
 }{2\delta_j}.
 \label{eq:app_central_difference}
\end{equation}
A scale-aware step may be defined by
\begin{equation}
 \delta_j
 =
 \epsilon^{1/3}
 \max\{1,|\theta_j|\}
 s_j,
 \label{eq:app_fd_step}
\end{equation}
where $\epsilon$ is machine precision and $s_j$ is an optional user-supplied
scale. One-sided differences are used near boundaries. The implementation
should support derivative checks based on alternative step sizes and report
relative discrepancies.

\subsubsection{Regularization}

If $\widehat\Omega_{gg}$, $\widehat S$, or $\widehat J$ is ill-conditioned,
a regularized matrix may be used:
\begin{equation}
 A_{\lambda}
 =
 A+\lambda I,
 \qquad \lambda\geq0.
 \label{eq:app_ridge_regularization}
\end{equation}
Alternatively, an eigenvalue floor may replace eigenvalues below
$\lambda_{\min}$. The regularization method, threshold, pre- and
post-regularization condition numbers, and effective numerical rank must be
reported. Regularization should be applied to the smallest set of matrices
needed to obtain a stable solution and should not be hidden inside generic
linear-algebra calls.

\subsubsection{Centering and scaling}

Moment centering must be identical across covariance blocks. Parameter and
moment scaling may improve numerical conditioning, but the transformation must
be invertible and recorded. Let $A_g$ and $A_h$ be nonsingular scaling
matrices. If transformed moments
\[
 g_i^*=A_gg_i,
 \qquad
 h_i^*=A_hh_i
\]
are used, the same transformations must be applied to covariance blocks,
Jacobians, projection coefficients, and diagnostics. Reported estimates and
standard errors are returned on the original parameter scale.

\subsubsection{Inner numerical solvers}

A demanding evaluation may itself contain a contraction mapping, simulation,
integration routine, or equilibrium solver. Each inner solver must return
\begin{enumerate}
 \item a convergence flag;
 \item the attained residual or tolerance;
 \item the number of iterations or function calls;
 \item whether a warm start was used; and
 \item any clipping, truncation, or fallback rule invoked.
\end{enumerate}
Moment evaluations based on unconverged inner solutions must be marked invalid.
The package may permit user-authorized fallback rules, but such rules must be
visible in the result object.

\subsubsection{Required statistical diagnostics}

The result object must report at least
\begin{align*}
 &\|\overline g(\widetilde\theta)\|,
 \qquad
 \|\overline h(\widetilde\theta)\|,
 \qquad
 \|\overline\nu\|,\\
 &\|\widehat\Delta_{g\nu}\|,
 \qquad
 \|\widehat d\|,
 \qquad
 \frac{\|\widehat d\|}{1+\|\widetilde\theta\|},\\
 &\kappa(\widehat\Omega_{gg}),
 \qquad
 \kappa(\widehat S),
 \qquad
 \kappa(\widehat J),
\end{align*}
where $\kappa(\cdot)$ denotes a reported condition-number convention. It should
also report effective ranks, regularization levels, and whether the final update
satisfies parameter constraints.

\subsubsection{Required computational diagnostics}

Evaluation accounting must distinguish
\begin{enumerate}
 \item tractable objective evaluations;
 \item tractable moment evaluations;
 \item demanding moment evaluations used for projection;
 \item demanding evaluations used for derivatives;
 \item tractable and demanding Jacobian evaluations;
 \item failed or rejected evaluations;
 \item inner-solver calls and iterations;
 \item final reconstruction evaluations; and
 \item elapsed time by computational stage.
\end{enumerate}
Let $N_H^{\mathrm{proj}}$, $N_H^{\mathrm{deriv}}$, and
$N_H^{\mathrm{recon}}$ denote the demanding evaluations used for projection,
derivatives, and reconstruction. The total is
\begin{equation}
 N_H^{\mathrm{total}}
 =
 N_H^{\mathrm{proj}}
 +N_H^{\mathrm{deriv}}
 +N_H^{\mathrm{recon}}.
 \label{eq:app_total_demanding_evals}
\end{equation}
These components must remain separately available even when a summary table
reports only their total.

\subsubsection{Warnings and validity checks}

An implementation should issue explicit warnings when
\begin{enumerate}
 \item localization fails or terminates on a parameter boundary;
 \item any moment or derivative contains a non-finite value;
 \item a demanding inner solver fails to converge;
 \item covariance blocks have deficient numerical rank;
 \item the orthogonality residual exceeds a user-defined tolerance;
 \item the one-step correction is large relative to the preliminary estimate;
 \item substantial regularization is required; or
 \item inference is requested under a covariance convention inconsistent with
 the supplied elementary unit.
\end{enumerate}
Warnings should not automatically invalidate an estimate, but they must remain
attached to the returned result.

%%%%%%%%%%%%%%%%%%%%%%%%%%%%%%%%%%%%%%%%%%%%%%%%%%%%%%%%%%%%%%%%%%%%%%%%%%%
\subsection{Software Architecture}
\label{app:software_architecture}
%%%%%%%%%%%%%%%%%%%%%%%%%%%%%%%%%%%%%%%%%%%%%%%%%%%%%%%%%%%%%%%%%%%%%%%%%%%

The specification implies a modular software design with four layers.

\paragraph{Model layer.}
The model layer owns the data, parameter transformation, moment operators,
derivative routines, reconstruction routines, and model-specific numerical
state. It does not construct the generic projection estimator.

\paragraph{Estimation layer.}
The estimation layer performs tractable localization, requests local demanding
evaluations, constructs projection objects, computes the one-step update, and
forms inference. It depends only on the operator contract and is therefore
reusable across models.

\paragraph{Numerical layer.}
The numerical layer provides optimization, differentiation, linear algebra,
regularization, parallel execution, caching, and reproducible random-number
management. Numerical services should not alter the mathematical definition of
the moment operators.

\paragraph{Results layer.}
The results layer stores estimates, covariance matrices, moment values,
projection objects, Jacobians, numerical state, operation counts, timing,
warnings, and reconstructed economic quantities. A result must contain enough
information to reproduce the one-step update without rerunning the global
localization stage.

\paragraph{Minimal model contract.}
In software-neutral notation, a compatible model supplies
\begin{verbatim}
tractable_moments(theta)   -> unit-by-q_g array
 demanding_moments(theta) -> unit-by-q_h array
tractable_jacobian(theta)  -> q_g-by-p array
 demanding_jacobian(theta)-> q_h-by-p array
reconstruct(theta)         -> model-specific output
\end{verbatim}
together with unit identifiers, parameter bounds or transformations, and a
numerical-state object. Analytical Jacobian methods may be omitted when the
estimation layer is authorized to construct numerical derivatives.

\paragraph{Caching and state consistency.}
A cache key must include the parameter vector, fidelity level, numerical state,
and any tolerance affecting the returned value. Reusing a cached demanding
evaluation under a different simulation draw set, contraction tolerance, or
fidelity level is invalid. Cache hits and misses should be counted because
caching can materially affect runtime comparisons.

\paragraph{Extensibility.}
The generic estimator must not contain model-specific references to BLP,
stochastic frontier models, mixed logit, or any other application. Extensions
are introduced by implementing the operator contract. This makes the package a
realization of the computational specification rather than a collection of
unrelated model-specific estimators.

%%%%%%%%%%%%%%%%%%%%%%%%%%%%%%%%%%%%%%%%%%%%%%%%%%%%%%%%%%%%%%%%%%%%%%%%%%%
\subsection{Specialization to Random-Coefficient Demand}
\label{app:blp_specialization}
%%%%%%%%%%%%%%%%%%%%%%%%%%%%%%%%%%%%%%%%%%%%%%%%%%%%%%%%%%%%%%%%%%%%%%%%%%%

The BLP application instantiates the computational partition through
low- and high-fidelity versions of the same structural demand operator. Let
$m_t^{L}(\theta)$ and $m_t^{H}(\theta)$ denote market-level moment contributions
constructed using low- and high-fidelity numerical settings. The elementary
statistical unit is the market $t=1,\ldots,T$. The computational partition is
\begin{equation}
 g_t(\theta)
 =
 m_t^{L}(\theta),
 \qquad
 h_t(\theta)
 =
 m_t^{H}(\theta)-m_t^{L}(\theta).
 \label{eq:app_blp_partition}
\end{equation}
Thus the complete high-fidelity moment is recovered as
\begin{equation}
 m_t^{H}(\theta)
 =
 g_t(\theta)+h_t(\theta).
 \label{eq:app_blp_recovery}
\end{equation}
The low-fidelity operator is a localization device, not an alternative economic
model. Both fidelity levels use the same structural parameterization,
instruments, product data, market definition, and economic target. They differ
only in numerical accuracy or computational intensity.

\subsubsection{BLP operator contract}

The model-specific implementation supplies the following operations.

\paragraph{Share inversion.}
Given $\theta$ and a fidelity level $f\in\{L,H\}$, the inversion operator
returns mean utilities
\begin{equation}
 \delta^{f}(\theta)
 =
 \mathcal I(\theta;f),
 \label{eq:app_blp_inversion}
\end{equation}
together with contraction residuals, iteration counts, convergence status, and
warm-start information.

\paragraph{Moment construction.}
The moment operator maps the inverted mean utilities into market-level moments,
\begin{equation}
 m_t^{f}(\theta)
 =
 \mathcal M_t(\delta^{f}(\theta),\theta;f).
 \label{eq:app_blp_moments}
\end{equation}
The returned object must preserve market-level contributions rather than only
the aggregate instrument--error product.

\paragraph{Derivative construction.}
The derivative operator returns
\begin{equation}
 G(\theta)
 =
 \frac{\partial \overline m^{L}(\theta)}{\partial\theta'},
 \qquad
 H(\theta)
 =
 \frac{\partial
 [\overline m^{H}(\theta)-\overline m^{L}(\theta)]
 }{\partial\theta'}.
 \label{eq:app_blp_derivatives}
\end{equation}
Analytical derivatives, implicit derivatives of the share inversion, automatic
differentiation, or common-random-number finite differences may be used. The
method and fidelity-specific numerical settings must be reported.

\paragraph{Reconstruction.}
At the final estimate, a high-fidelity reconstruction may compute mean utilities,
linear parameters, marginal costs, elasticities, diversion ratios,
markups, or counterfactual outcomes. These calls are recorded as reconstruction
rather than projection evaluations.

\subsubsection{Fidelity specification}

A fidelity level is a complete numerical configuration. It may differ in
\begin{enumerate}
 \item the number of simulation draws;
 \item integration or quadrature rules;
 \item contraction tolerance and maximum iterations;
 \item acceleration method;
 \item derivative accuracy;
 \item fixed-point initialization; or
 \item the accuracy of an inner supply-side or equilibrium calculation.
\end{enumerate}
For a valid comparison, all features not intended to define fidelity must be
held fixed. In particular, the same data, instruments, objective scaling,
parameter transformations, and random-number design should be used across the
two operators.

\subsubsection{BLP evaluation accounting}

For each fidelity level, the implementation records
\begin{enumerate}
 \item calls to simulated-share evaluation;
 \item contraction calls and total contraction iterations;
 \item Jacobian calls;
 \item objective and gradient evaluations;
 \item failed contractions or clipped shares;
 \item simulation draws and random seed;
 \item time spent in simulation, contraction, differentiation, linear
 concentration, and matrix operations; and
 \item final high-fidelity reconstruction calls.
\end{enumerate}
Let $N_H^{\mathrm{GMM}}$ denote the number of high-fidelity evaluations used by
conventional full GMM and $N_H^{\mathrm{SEIP}}$ the corresponding number used by
the projected estimator. The structural evaluation ratio is
\begin{equation}
 \mathcal R_H
 =
 \frac{N_H^{\mathrm{SEIP}}}
 {N_H^{\mathrm{GMM}}}.
 \label{eq:app_blp_evaluation_ratio}
\end{equation}
Runtime comparisons should be accompanied by $\mathcal R_H$ and the underlying
operation counts. This separates the structural computational gain from
hardware, parallelization, and implementation overhead.

\subsubsection{BLP validity diagnostics}

The BLP specialization should report
\begin{enumerate}
 \item the high- and low-fidelity moment difference at
 $\widetilde\theta$;
 \item the norm of the projected high-fidelity correction;
 \item derivative agreement across alternative local step sizes;
 \item contraction convergence at every demanding evaluation;
 \item sensitivity to the low- and high-fidelity definitions;
 \item parameter and standard-error differences relative to full
 high-fidelity GMM; and
 \item differences in reconstructed economic quantities.
\end{enumerate}
A large, unstable, or fidelity-sensitive correction is evidence against the
chosen low-fidelity operator as an adequate localizer. The package should expose
this evidence rather than treating successful numerical termination as
automatic validation of the approximation.

%%%%%%%%%%%%%%%%%%%%%%%%%%%%%%%%%%%%%%%%%%%%%%%%%%%%%%%%%%%%%%%%%%%%%%%%%%%
\subsection{Mapping to the \texttt{OrthoGMM} Package}
\label{app:orthogmm_mapping}
%%%%%%%%%%%%%%%%%%%%%%%%%%%%%%%%%%%%%%%%%%%%%%%%%%%%%%%%%%%%%%%%%%%%%%%%%%%

The package should preserve a one-to-one mapping between the mathematical
objects in this appendix and stored software objects. At minimum, the mapping is
\begin{center}
\begin{tabular}{ll}
\hline
Mathematical object & Package role \\
\hline
$g_i(\theta)$ & unit-level tractable moments \\
$h_i(\theta)$ & unit-level demanding moments \\
$\mathcal G$ & tractable moment operator \\
$\mathcal H$ & demanding moment operator \\
$\mathcal D$ & derivative provider \\
$\mathcal R$ & reconstruction provider \\
$\widetilde\theta$ & tractable localizer \\
$\widehat\Omega_{gg},\widehat\Omega_{hg}$ & covariance blocks \\
$\widehat B$ & projection coefficient \\
$\widehat\nu_i$ & orthogonal demanding residuals \\
$\widehat S$ & residual covariance \\
$\widehat G,\widehat H,\widehat R$ & Jacobian objects \\
$\widehat J$ & projected information matrix \\
$\widehat d$ & efficient one-step correction \\
$\widehat\theta_{\mathrm{SEIP}}$ & final estimator \\
$n^{-1}\widehat J^{-1}$ & baseline covariance estimate \\
\hline
\end{tabular}
\end{center}

A serialized result should contain these objects whenever their dimensions make
storage practical. For large unit-level arrays, the result may instead retain a
reproducible reference to cached data together with checksums or equivalent
integrity information.

The package should expose separate estimator classes or modes for tractable GMM,
full GMM, and projected GMM, but all should share the same model object and
numerical-state conventions. This permits direct comparisons in which data,
instruments, moment scaling, parameter transformations, and simulation draws are
held fixed. Reported differences can then be attributed to the computational
route rather than inconsistent model implementations.

Finally, the package documentation and empirical scripts should identify this
appendix as the authoritative estimator specification. Software behavior that
departs from the equations or workflow stated here should be documented as an
extension, approximation, or experimental option rather than silently treated
as the canonical estimator.
